\section{The zero-Weyl endpoint at finite Euler density}
\label{sec:finite-density-endpoint}

The endpoint estimate differs from the finite-parameter inequality in
Section~\ref{sec:single-parameter}: its deficit is $D_\beta/V$, not
$D_\beta/(\beta V)$.  A bound on Euler characteristic relative to the
positive projection budget suffices to remove the decay at
$\beta=\infty$.  The proof uses the actual monopole to construct a
section of the fifth power of its determinant line.  A logarithmically
selected resolvent then detects the nonzero characteristic square of
that line, even when the underlying manifolds degenerate.

\begin{theorem}[Finite-density endpoint]
\label{thm:finite-density-endpoint}
For every finite $B\ge0$ there is a constant $c(B)>0$ with the following
property.  Let $M$ be a closed connected oriented smooth four-manifold,
and assume the actual conformal monopole hypothesis of
Section~\ref{sec:single-parameter}: a fixed
$\operatorname{Spin}^c$ structure admits an ordinary unperturbed
monopole solution for every metric in the conformal class under
consideration.  Let $a=c_1(L)$ be its determinant class and suppose
\[
 p=(a_g^+)^2\ge A>0,\qquad a^2\ge T_+=2\chi(M)+3\tau(M),
 \qquad \frac{\chi(M)}{A}\le B.
\]
Then every metric $g$ in that conformal class and every finite
$\beta\ge2$ satisfy
\begin{equation}
 \kappa_\beta(g):=
 \fint_M|\Ric_g+g|^2+\frac{D_\beta(g)}{V(g)}
 \ge c(B)\frac{(\beta-2)^2}{\beta^2},
 \label{eq:finite-density-gap}
\end{equation}
where
\[
 D_\beta(g)=\left(1+\frac1\beta\right)\int_Ms_g^2
       +6(\beta+1)\int_M|W_g^+|^2-72\pi^2p.
\]
The constant is independent of the manifold, the selected structure,
$A$, the metric, and $\beta$ within the displayed hypotheses.  No
pointwise Ricci bound, diameter bound, noncollapse, or uniform Sobolev
inequality is assumed.
\end{theorem}

The structure may be chosen separately on different conformal classes;
the same structure is retained throughout each class.  On a
representative with $s=-4$, the proof uses only one ordinary solution,
rescaled to the normalization below.  Thus the original conformal
supplier is used to reach that representative, not to choose a family
of endpoint solutions as $\beta$ varies.  Throughout this section,
$\|\cdot\|_p$ and $\fint$ denote norms and integrals with respect to
$d\mu_g/V(g)$; $\|\cdot\|_\infty$ is the pointwise essential supremum.
For a four-form, $\fint$ means its integral divided by $V(g)$.

\subsection{The actual strong-current input}

We first state precisely the upstream input used by the endpoint
argument.  Its strong critical estimates are proved in
Section~\ref{sec:web-strong-current}; the constant-weight identities
and the bounded anti-self-dual curvature fit are in
Section~\ref{sec:single-parameter}.

\begin{lemma}[Strong current and bounded determinant curvature]
\label{lem:finite-density-strong-current}
Suppose that a sequence satisfying the actual supply and square
hypotheses of Theorem~\ref{thm:finite-density-endpoint} has
\begin{equation}
 s_{g_j}=-4,\qquad \beta_j\longrightarrow\infty,\qquad
 \kappa_j=\frac{P_j+D_{\beta_j}}{V_j}\longrightarrow0.
 \label{eq:finite-density-zero-sequence}
\end{equation}
An Euler-density bound is not required for the conclusions of this
lemma.  Choose one actual solution on each member, with determinant
connection $\mathcal A_j$, so that
\begin{equation}
 D_{\mathcal A_j}^+\Phi_j=0,\qquad
 F_{\mathcal A_j}^+=2i\sigma(\Phi_j),\qquad
 \rho_j=|\Phi_j|^2\le2,
 \qquad
 \nabla_{\mathcal A_j}^*\nabla_{\mathcal A_j}\Phi_j
       =(1-\rho_j/2)\Phi_j.
 \label{eq:finite-density-actual-equations}
\end{equation}
Define the actual cubic torsion by
\begin{equation}
 t_j=t(\Phi_j):=\sum_i c(e^i)\Phi_j\otimes
       (\Phi_j\wedge\nabla_{\mathcal A_j,e_i}\Phi_j)
       \in\Gamma(E_{-,j}),\qquad E_{-,j}=S_j^-\otimes L_j.
 \label{eq:finite-density-cubic-torsion}
\end{equation}
Its determinant line is $L_j^3$; write $D_j=D^-_{3\mathcal A_j}$
and $\mathcal L_j=\nabla_{3\mathcal A_j}^*\nabla_{3\mathcal A_j}$.
With
\[
 \rho_*=\frac43,\qquad m_*=\frac{128}{81},\qquad m_j=|t_j|^2,
\]
one has
\begin{gather}
 \rho_j\longrightarrow\rho_* \text{ in every finite }L^p,
 \qquad m_j\longrightarrow m_* \text{ in }L^2,
 \label{eq:finite-density-strong-masses}\\
 \|\nabla\Phi_j\|_4+\|\nabla^2\Phi_j\|_2+
 \|t_j\|_4+\|\nabla t_j\|_2\le C,
 \qquad \|D_jt_j\|_2\longrightarrow0.
 \label{eq:finite-density-strong-bounds}
\end{gather}
On every fixed positive-$\rho$ region the full $(1,0)$ projection
of $\nabla t_j$ for the almost-complex structure $J_{+,j}$ selected
by $\Phi_j$ tends to zero in $L^2$.  Moreover
\begin{equation}
 (\mathcal L_j-2)t_j
   =D_j^*D_jt_j+(\Gamma_j-1)t_j-\mathcal R_jt_j,
 \quad
 |\Gamma_j|\le C,\quad
 \Gamma_j\longrightarrow1 \text{ in measure},\quad
 \|\mathcal R_j\|_2\longrightarrow0.
 \label{eq:finite-density-torsion-equation}
\end{equation}
Here $\mathcal R_j$ is an endomorphism of $E_{-,j}$.  The full real
determinant curvature admits a decomposition
\begin{equation}
 -iF_{\mathcal A_j}=F_{0,j}+r_{F,j},\qquad
 |F_{0,j}|\le C,\quad \|r_{F,j}\|_2\longrightarrow0,
 \label{eq:finite-density-determinant-fit}
\end{equation}
and the characteristic numbers satisfy
\begin{equation}
 \frac{a_j^2}{V_j}\longrightarrow\frac1{6\pi^2}.
 \label{eq:finite-density-characteristic-limit}
\end{equation}
\end{lemma}

We recall the ingredients of this input to specify the strength of the
convergence being used.  Put $\zeta=\langle\Phi,\nabla\Phi\rangle$
and $x=\kappa_\beta+\beta^{-2}$.  For a fixed $\epsilon_0>0$, the
derivative-free weighted Hessian calculation in
Section~\ref{sec:web-strong-current} gives, for every $L_0\ge1$,
\begin{equation}
 \fint
 \left(\frac{|\nabla\Phi|^2}{L_0^2+|\nabla\Phi|^2}\right)^{\epsilon_0}
 \bigl(|\nabla^2\Phi|^2+|\nabla\Phi|^4\bigr)
 \le C L_0^{-2\epsilon_0}+C\sqrt{x}.
 \label{eq:finite-density-slow-truncation}
\end{equation}
The positive-spinor curvature has a bounded model and an $L^2$-small
remainder; it contains no $W^-$ block.  Consequently
\eqref{eq:finite-density-slow-truncation} makes
$|\nabla\Phi|^4$ uniformly integrable on the normalized measure
spaces.  Splitting at $|\nabla\Phi|=L_0$, first letting $j\to\infty$
and then $L_0\to\infty$, proves
\[
 \fint |\zeta|^2|\nabla\Phi|^2\longrightarrow0.
\]
The exact cubic Dirac identity and the current identities then give
$\|Dt\|_2\to0$, $\||t|^2-m_*\|_2\to0$, and
$\|\nabla\zeta\|_2\to0$.  The last assertion supplies the full
first projection, rather than only its Dirac trace.  These steps,
including removal of the low-$\rho$ sets, are part of the proof in
Section~\ref{sec:web-strong-current}.

For completeness, the same-field anti-self-dual model used there is
\[
 Q=-iq_-(t),\qquad
 S_0=m_*+2\rho^2|\nabla\Phi|^2,\qquad
 K=-\frac98\frac{\rho^3Q}{S_0}.
\]
It satisfies $|K|\le C$ and
\begin{equation}
 \|-iF_{\mathcal A}^--K\|_2^2
       +\frac{4\pi^2(a^2-T_+)}V\le C\sqrt{x}.
 \label{eq:finite-density-asd-fit}
\end{equation}
Together with the bounded self-dual part in
\eqref{eq:finite-density-actual-equations}, this proves
\eqref{eq:finite-density-determinant-fit}.  In
\eqref{eq:finite-density-torsion-equation} one may take
$\Gamma=(27/32)\rho^3|t|^2/S_0$.  Finally $W^+\to0$ in $L^2$ and
the chiral Chern--Weil formula give $T_+/V\to1/(6\pi^2)$;
\eqref{eq:finite-density-asd-fit} gives the same limit for $a^2/V$.
This does not assert uniform integrability of $|\nabla t|^4$ or of
$|\nabla^2\Phi|^2$ on arbitrary small sets.

\subsection{A rank-one resolvent obstruction}

The following analytic lemma is independent of monopole theory.  The
closed testing form need not have a closed bounded model.

\begin{lemma}[Logarithmic resolvent and line-bundle transgression]
\label{lem:rank-one-logarithmic-resolvent}
Let $(M_j,g_j)$ be closed oriented Riemannian four-manifolds, and let
$K_j$ be Hermitian line bundles with unitary connections $\nabla_j$.
Write $F_j=-iF_{\nabla_j}$ for their real curvature and assume
\[
 \sup_j\bigl(\|\Ric_{g_j}\|_2^2+\|F_j\|_2^2\bigr)<\infty.
\]
Suppose smooth sections $\sigma_j$ satisfy
\[
 |\sigma_j|\longrightarrow\ell>0 \text{ in normalized measure},
 \qquad \sup_j\|\nabla_j\sigma_j\|_2<\infty.
\]
If real closed two-forms $\omega_j$ satisfy
\[
 \omega_j=\omega_{0,j}+e_j,\qquad
 \|\omega_{0,j}\|_\infty\le C,\qquad \|e_j\|_2\longrightarrow0,
\]
then
\begin{equation}
 \frac1{V_j}\int_{M_j}F_j\wedge\omega_j\longrightarrow0.
 \label{eq:finite-density-line-pairing}
\end{equation}
In particular, if $F_j$ itself has this bounded-model decomposition,
then
\begin{equation}
 \frac{c_1(K_j)^2[M_j]}{V_j}\longrightarrow0.
 \label{eq:finite-density-line-square}
\end{equation}
\end{lemma}

\begin{proof}
Choose a smooth unitary-equivariant radial map $T$ on the line,
equal to the identity near $|z|=\ell$, with bounded image and bounded
derivative.  The sections $s_j=T(\sigma_j)$ satisfy
\[
 |s_j|\le M_0,\qquad \|\nabla_js_j\|_2^2\le E_0,
 \qquad d_j:=\||s_j|-\ell\|_2^2\longrightarrow0.
\]
Boundedness and convergence in measure justify the last limit without
an integrability assumption on the original sections.

Let $\mathcal L_j=\nabla_j^*\nabla_j$.  For $\varepsilon>0$ define
$u_{j,\varepsilon}=(I+\varepsilon\mathcal L_j)^{-1}s_j$.
The connection-Laplacian maximum principle and spectral calculus give
\begin{gather}
 |u_{j,\varepsilon}|\le M_0,\qquad
 \|u_{j,\varepsilon}-s_j\|_2^2\le\varepsilon E_0,
 \qquad\|\nabla_ju_{j,\varepsilon}\|_2^2\le E_0,
 \label{eq:finite-density-line-resolvent}\\
 \int_0^\infty\|\mathcal L_ju_{j,\varepsilon}\|_2^2\,d\varepsilon
       =\|\nabla_js_j\|_2^2\le E_0.
 \label{eq:finite-density-log-average}
\end{gather}
For each eigenvalue $\lambda\ge0$, the identity in the second line is
$\int_0^\infty\lambda^2(1+\varepsilon\lambda)^{-2}
\,d\varepsilon=\lambda$.  No uniform spectral gap is used.

Set $\delta_j=d_j+\|e_j\|_2^2+j^{-1}\to0$.  Averaging
\eqref{eq:finite-density-log-average} with respect to
$d\varepsilon/\varepsilon$ on
$[\delta_j^{1/2},\delta_j^{1/4}]$ gives a choice
$\varepsilon_j$ in that interval for which, on writing
$u_j=u_{j,\varepsilon_j}$ and
$k_j=\|\mathcal L_ju_j\|_2^2$,
\begin{equation}
 \varepsilon_jk_j\le\frac{4E_0}{|\log\delta_j|}\longrightarrow0,
 \qquad
 \frac{\|e_j\|_2^2}{\varepsilon_j}\le\delta_j^{1/2}
       \longrightarrow0.
 \label{eq:finite-density-log-selection}
\end{equation}
An arbitrary smoothing parameter would yield only
$\varepsilon_jk_j=O(1)$, which is insufficient below.

We record the derivative estimate used for these resolvents.  For a
bounded smooth section $u$ of any fixed-rank unitary bundle, put
$I=\fint|\nabla^2u|^2$ and $J=\fint|\nabla u|^4$.
Integrating the divergence of
$|\nabla u|^2\operatorname{Re}\langle u,\nabla u\rangle$ gives
\[
 J\le C\|u\|_\infty I^{1/2}J^{1/2},\qquad
 J\le C\|u\|_\infty^2 I.
\]
Apply the integrated covariant Hodge--Weitzenb\"ock identity to the
bundle-valued one-form $\nabla u$.  Its covariant exterior derivative
is $F_\nabla u$ and its codifferential is $\nabla^*\nabla u$;
the curvature on the covector slot contracts to Ricci.  Hence
\[
 I\le\|\nabla^*\nabla u\|_2^2
   +C\fint(|\Ric|+|F_\nabla|)|\nabla u|^2
   +C\|u\|_\infty^2\|F_\nabla\|_2^2.
\]
Cauchy--Schwarz and Young's inequality, using the preceding
interpolation, imply
\begin{equation}
 \fint\bigl(|\nabla^2u|^2+|\nabla u|^4\bigr)
 \le C_{M_0}\bigl(
 \|\nabla^*\nabla u\|_2^2+
 \|\Ric\|_2^2+\|F_\nabla\|_2^2\bigr),
 \qquad |u|\le M_0.
 \label{eq:finite-density-hessian-interchange}
\end{equation}
This uses no curvature derivative or Sobolev inequality.  In
particular the fourth-power derivative integral of $u_j$ is at most
$C(1+k_j)$.

Fix $0<\ell_0<\ell_1<\ell$ and put
$Z_j=\{|u_j|\le\ell_1\}$.  Chebyshev's inequality and
\eqref{eq:finite-density-line-resolvent} give
\[
 \mu_j(Z_j)\le C(d_j+E_0\varepsilon_j)\le C\varepsilon_j,
 \qquad \mu_j=\frac{d\mu_{g_j}}{V_j}.
\]
Consequently
\begin{align}
 \fint_{Z_j}|\nabla u_j|^2
 &\le C\sqrt{\varepsilon_j(1+k_j)}\longrightarrow0,
 \label{eq:finite-density-transition-energy}\\
 \|e_j\|_2\left(\fint|\nabla u_j|^4\right)^{1/2}
 &\le C\sqrt{\|e_j\|_2^2(1+k_j)}\longrightarrow0.
 \label{eq:finite-density-transition-error}
\end{align}
For the second limit, use
$\|e_j\|_2^2k_j=(\|e_j\|_2^2/\varepsilon_j)
(\varepsilon_jk_j)$ and \eqref{eq:finite-density-log-selection}.

On $\{u_j\ne0\}$ the unit section $u_j/|u_j|$ has a real connection
form $\alpha_j$ with
\[
 d\alpha_j=F_j,\qquad
 |\alpha_j|\le C|\nabla u_j|/|u_j|.
\]
Choose a fixed smooth function $\chi$ which is zero on
$[0,\ell_0]$ and one on $[\ell_1,\infty)$, and set
$\chi_j=\chi(|u_j|)$.  Then $\chi_j\alpha_j$ extends smoothly by
zero to all of $M_j$.  Since $d\omega_j=0$, Stokes' theorem gives
\begin{equation}
 \fint F_j\wedge\omega_j
 =\fint(1-\chi_j)F_j\wedge\omega_j
       -\fint d\chi_j\wedge\alpha_j\wedge\omega_j.
 \label{eq:finite-density-line-transgression}
\end{equation}
The first term is bounded in absolute value by
\[
 C\|F_j\|_2\mu_j(Z_j)^{1/2}
       +\|F_j\|_2\|e_j\|_2\longrightarrow0.
\]
On the support of $d\chi_j$,
$|d\chi_j\wedge\alpha_j|\le C|\nabla u_j|^2$.  Splitting
$\omega_j=\omega_{0,j}+e_j$ bounds the second term by
\[
 C\fint_{Z_j}|\nabla u_j|^2
 +C\|e_j\|_2\left(\fint|\nabla u_j|^4\right)^{1/2}
       \longrightarrow0
\]
by \eqref{eq:finite-density-transition-energy}--
\eqref{eq:finite-density-transition-error}.  This proves
\eqref{eq:finite-density-line-pairing}.  For $\omega_j=F_j$, the
identity $c_1(K_j)=[F_j/(2\pi)]$ gives
\eqref{eq:finite-density-line-square}.
\end{proof}

\subsection{A bounded opposite-spinor resolvent}

Assume now that a sequence as in
\eqref{eq:finite-density-zero-sequence} also satisfies
$\chi(M_j)/A_j\le B<\infty$.  We shall derive a contradiction.
The scalar monopole inequality gives
\[
 16V_j\ge32\pi^2p_j\ge32\pi^2A_j.
\]
The Euler identity therefore supplies the missing opposite-Weyl bound:
\begin{equation}
 b_j:=\fint|W_j^-|^2
 =8\pi^2\frac{\chi(M_j)}{V_j}
       -\fint|W_j^+|^2-\frac23+\frac{P_j}{2V_j}
 \le4B+\frac{P_j}{2V_j}.
 \label{eq:finite-density-opposite-weyl}
\end{equation}
All constants depending on $B$ below arise through this bound.

Choose a smooth radial cap $T$ on $E_-$, equal to the identity for
$|t|\le2$ and bounded by $4$, and let $z_j=T(t_j)$.  Since
$m_*<4$ and $|t_j|^2\to m_*$ in $L^2$,
$\|z_j-t_j\|_2\to0$.  Define the exact resolvent
\begin{equation}
 (4+\mathcal L_j)v_j=6z_j.
 \label{eq:finite-density-direct-resolvent}
\end{equation}
Its maximum principle gives $|v_j|\le6$ and
$|\mathcal L_jv_j|\le48$.  To verify the required strong
approximation, set $w_j=v_j-t_j$ and $d_j=D_jt_j$.  Equation
\eqref{eq:finite-density-torsion-equation} gives
\[
 (4+\mathcal L_j)w_j
 =6(z_j-t_j)-D_j^*d_j-(\Gamma_j-1)t_j+\mathcal R_jt_j.
\]
Testing against $w_j$, absorbing the first two terms by Young's
inequality and $|D_jw_j|\le2|\nabla w_j|$, gives
\begin{align*}
 4\|w_j\|_2^2+\|\nabla w_j\|_2^2
 ={}&6\operatorname{Re}\fint\langle z_j-t_j,w_j\rangle
       -\operatorname{Re}\fint\langle d_j,D_jw_j\rangle\\
 &-\operatorname{Re}\fint\langle(\Gamma_j-1)t_j,w_j\rangle
       +\operatorname{Re}\fint\langle\mathcal R_jt_j,w_j\rangle.
\end{align*}
The last two terms tend to zero because
\[
 \bigl|\fint\langle(\Gamma_j-1)t_j,w_j\rangle\bigr|
 +\bigl|\fint\langle\mathcal R_jt_j,w_j\rangle\bigr|
 \le (\|\Gamma_j-1\|_2+\|\mathcal R_j\|_2)
       \|t_j\|_4\|w_j\|_4=o(1).
\]
Here $\|w_j\|_4\le\|v_j\|_\infty+\|t_j\|_4\le C$ directly on
the normalized probability space; no Sobolev embedding is invoked.
Thus
\begin{equation}
 \|v_j-t_j\|_{H^1}\longrightarrow0,\qquad
 e_j:=(\mathcal L_j-2)v_j=6(z_j-v_j)\longrightarrow0
       \text{ in }L^2,\qquad
 \|D_jv_j\|_2\longrightarrow0.
 \label{eq:finite-density-resolvent-convergence}
\end{equation}
This is the direct construction of
Section~\ref{sec:web-direct-resolvent}; it does not differentiate the
cap $T(t_j)$.

The curvature of $E_-$ contains $W^-$, the Ricci blocks, the scalar
block, and the determinant curvature.  By
\eqref{eq:finite-density-opposite-weyl} and
Lemma~\ref{lem:finite-density-strong-current}, its normalized $L^2$
norm is bounded by $C_B$.  Applying
\eqref{eq:finite-density-hessian-interchange} to the bounded $v_j$
therefore yields
\begin{equation}
 \|\nabla^2v_j\|_2+\|\nabla v_j\|_4^2\le C_B.
 \label{eq:finite-density-negative-hessian}
\end{equation}
Put $\alpha_j=|v_j|^2$ and $q_j=|\nabla v_j|^2$.  Boundedness of
$v_j$ and its convergence in measure to the amplitude of $t_j$ imply
$\alpha_j\to m_*$ in every finite $L^p$.  With the nonnegative
scalar Laplacian, its exact scalar equation is
\begin{equation}
 \Delta\alpha_j=4\alpha_j-2q_j
       +2\operatorname{Re}\langle e_j,v_j\rangle.
 \label{eq:finite-density-scalar-mass}
\end{equation}
The right side is bounded in $L^2$ by
\eqref{eq:finite-density-negative-hessian}; testing with
$\alpha_j-m_*$ gives
\begin{equation}
 \|d\alpha_j\|_2^2
 \le\|\alpha_j-m_*\|_2\|\Delta\alpha_j\|_2\longrightarrow0.
 \label{eq:finite-density-mass-gradient}
\end{equation}

\subsection{Mass, phase, and the surviving derivative}

For clarity we suppress $j$ in the frame calculations.  Work on fixed
positive-amplitude sets for $\Phi$ and $v$.  Choose compatible
orthonormal tangent and unitary spinor frames with
$\Phi=\sqrt\rho\,(1,0)$ and $v=\sqrt\alpha\,(1,0)$, and put
\[
 Z_1=e_1-ie_2,\qquad Z_2=e_3-ie_4,
 \qquad j_v=\operatorname{Im}\langle v,\nabla v\rangle.
\]
The full $(1,0)$ projection $N_v$ tends to zero in $L^2$ by the
strong-current lemma and \eqref{eq:finite-density-resolvent-convergence}.
The exact mass--phase identity is
\begin{equation}
 |j_v+\tfrac12J_+d\alpha|^2\le2\alpha|N_v|^2.
 \label{eq:finite-density-phase}
\end{equation}
Hence $j_v\to0$ in $L^2$ on those sets.  Each first spinor
component obeys, first for real $X$ and then by complex linearity,
\[
 (\nabla_Xv)_1
       =\frac{d\alpha(X)/2+ij_v(X)}{\sqrt\alpha}.
\]
The Clifford formula in the compatible frame is, up to an overall
sign,
\[
 D^-v=\bigl((\nabla_{Z_1}v)_1-(\nabla_{Z_2}v)_2,
            (\nabla_{\bar Z_1}v)_2+(\nabla_{\bar Z_2}v)_1\bigr).
\]
Combining these identities proves
\begin{equation}
 \nabla_{Z_1}v\longrightarrow0,\quad
 \nabla_{Z_2}v\longrightarrow0,\quad
 \nabla_{\bar Z_1}v\longrightarrow0,\quad
 (\nabla_{\bar Z_2}v)_1\longrightarrow0
       \quad\text{in }L^2.
 \label{eq:finite-density-first-projections}
\end{equation}
The assertion is invariant under compatible frame changes and is
understood with fixed smooth amplitude cutoffs.  Indeed, for any fixed
$0<\rho_0<\rho_*$ and $0<\alpha_0<m_*$, the set
$E_j=\{\rho_j<\rho_0\}\cup\{\alpha_j<\alpha_0\}$ has
$\mu_j(E_j)\to0$.  Its full quadratic derivative contribution obeys
\[
 \fint_{E_j}|\nabla v_j|^2
       \le\mu_j(E_j)^{1/2}\|\nabla v_j\|_4^2\longrightarrow0.
\]
The phase has the same property since $v_j$ is bounded.  This
justifies removing the exceptional sets in
\eqref{eq:finite-density-first-projections}; neither a global frame
nor a lower pointwise amplitude bound has been assumed.

It remains to determine the modulus of the surviving derivative.
Let $q_*=2m_*$ and $T_0(q)=q/(1+q)$.  The Kato estimate gives
\[
 |dT_0(q_j)|\le\frac{2\sqrt{q_j}}{(1+q_j)^2}|\nabla^2v_j|
       \le C|\nabla^2v_j|.
\]
Thus $T_0(q_j)-T_0(q_*)$ is a bounded test with uniformly bounded
$H^1$ seminorm.  Testing \eqref{eq:finite-density-scalar-mass} gives
\begin{align}
 2\fint(q_j-q_*)\bigl[T_0(q_j)-T_0(q_*)\bigr]
 ={}&4\fint(\alpha_j-m_*)\bigl[T_0(q_j)-T_0(q_*)\bigr]\notag\\
 &+2\fint\operatorname{Re}\langle e_j,v_j\rangle
                  \bigl[T_0(q_j)-T_0(q_*)\bigr]\notag\\
 &-\fint\langle d\alpha_j,dT_0(q_j)\rangle
       \longrightarrow0.
 \label{eq:finite-density-covariance}
\end{align}
Since
\[
 (q-q_*)[T_0(q)-T_0(q_*)]
       =\frac{(q-q_*)^2}{(1+q)(1+q_*)},
\]
we obtain $q_j\to2m_*$ in measure.  The $L^2$ bound on $q_j$
then gives strong $L^p$ convergence for every $1\le p<2$.
It does not give strong $L^2$ convergence of $q_j$.

Write $C_j=(\nabla_{\bar Z_2}v_j)_2$.  The unnormalized complex
directions chosen above satisfy
\[
 |\nabla v_j|^2=\frac12\left(
 |\nabla_{Z_1}v_j|^2+|\nabla_{\bar Z_1}v_j|^2+
 |\nabla_{Z_2}v_j|^2+|\nabla_{\bar Z_2}v_j|^2\right).
\]
It follows from \eqref{eq:finite-density-first-projections} and
\eqref{eq:finite-density-covariance} that
\begin{equation}
 |C_j|^2\longrightarrow4m_*\text{ in measure and }L^1,
 \qquad
 |c_j|^2\longrightarrow1\text{ in measure},\qquad
 c_j=\frac{C_j}{2\sqrt{\alpha_j}}.
 \label{eq:finite-density-surviving-modulus}
\end{equation}
The coefficient $c_j$ is used only on fixed positive-amplitude sets;
its modulus is invariant, and no global choice of its phase is needed.
These are the mass--phase--Dirac and covariance arguments of
Section~\ref{sec:web-weak-bianchi}.

\subsection{The actual fifth-power line and the contradiction}

The cubic operation \eqref{eq:finite-density-cubic-torsion} is defined
for an arbitrary spinor, whether or not it solves a monopole equation.
For $v\in S^-\otimes L$, its determinant wedge has values in
$\det(S^-\otimes L)=L^3$.  Therefore
\begin{equation}
 t(v)=\sum_i c(e^i)v\otimes(v\wedge\nabla_{e_i}v)
       \in\Gamma(S^+\otimes L^4),\qquad
 \sigma=\Phi\wedge t(v)\in\Gamma(L^5).
 \label{eq:finite-density-fifth-line}
\end{equation}
Both sections are globally smooth polynomial expressions.  The
exponent $5$ uses $\det S^+=L$ in the last wedge; it is not a
choice of an auxiliary trivialization.

In the compatible frames above, the surviving first jet has
off-diagonal connection coefficient $(0,0,c,-ic)$ and gives
\[
 t(v)=(0,2\alpha^{3/2}c),\qquad
 |\Phi\wedge t(v)|^2=4\rho\alpha^3|c|^2.
\]
The terms omitted from this ideal-jet expression are bounded linear
multiples of the projections in
\eqref{eq:finite-density-first-projections}, since $\Phi$ and $v$
are bounded.  They tend to zero in $L^2$ on the positive-amplitude
sets; the exceptional sets have vanishing measure.  Hence
\begin{equation}
 |\sigma_j|^2\longrightarrow
       4\rho_*m_*^3>0\quad\text{in normalized measure}.
 \label{eq:finite-density-fifth-mass}
\end{equation}
The polynomial formula also gives the global estimates
\begin{gather*}
 |\sigma_j|\le C|\nabla v_j|,\\
 |\nabla_{5\mathcal A_j}\sigma_j|
       \le C\bigl(|\nabla\Phi_j|\,|\nabla v_j|
                    +|\nabla v_j|^2+|\nabla^2v_j|\bigr).
\end{gather*}
H\"older's inequality and
\eqref{eq:finite-density-strong-bounds},
\eqref{eq:finite-density-negative-hessian} therefore imply
\begin{equation}
 \|\sigma_j\|_4+\|\nabla_{5\mathcal A_j}\sigma_j\|_2\le C_B.
 \label{eq:finite-density-fifth-energy}
\end{equation}
In particular $|\sigma_j|$ converges strongly in $L^2$ to the
positive square root of the limit in
\eqref{eq:finite-density-fifth-mass}.  No nowhere-zero section has
been asserted.

Apply Lemma~\ref{lem:rank-one-logarithmic-resolvent} to $K_j=L_j^5$
and $\omega_j=-iF_{5\mathcal A_j}=5(-iF_{\mathcal A_j})$.
Its hypotheses hold: $\|\Ric_{g_j}\|_2^2=4+P_j/V_j$ is bounded,
the line curvature has the bounded-model decomposition
\eqref{eq:finite-density-determinant-fit}, and
\eqref{eq:finite-density-fifth-mass}--
\eqref{eq:finite-density-fifth-energy} supply the section estimates.
The lemma gives
\[
 \frac{c_1(L_j^5)^2[M_j]}{V_j}\longrightarrow0.
\]
On the other hand,
\eqref{eq:finite-density-characteristic-limit} gives
\begin{equation}
 \frac{c_1(L_j^5)^2[M_j]}{V_j}
       =25\frac{a_j^2}{V_j}
       \longrightarrow\frac{25}{6\pi^2}>0,
 \label{eq:finite-density-contradiction}
\end{equation}
a contradiction.  This excludes every sequence in
\eqref{eq:finite-density-zero-sequence} with fixed finite $B$.

\subsection{All parameters and the original metric}

\begin{proof}[Completion of the proof of
Theorem~\ref{thm:finite-density-endpoint}]
The conformal supply and $p>0$ force the Yamabe constant to be
negative, by the spinor Weitzenb\"ock identity on a constant-scalar
representative.  Let $\widehat g$ denote the representative with
$s_{\widehat g}=-4$.  The Hodge projection $p$, the class square,
and both integrated Weyl energies are conformally invariant.  If
$\delta=S(g)-S(\widehat g)$, the exact conformal identities in
Section~\ref{sec:single-parameter} give
\begin{gather}
 \delta\ge0,\qquad
 \widehat P=P-\delta/12,\qquad
 \widehat D_\beta=D_\beta-(1+\beta^{-1})\delta\ge0,
 \label{eq:finite-density-conformal-identities}\\
 \widehat\kappa_\beta
       :=\frac{\widehat P+\widehat D_\beta}{\widehat V}
       \le\frac{\kappa_\beta}
        {1-\sqrt{\kappa_\beta}-\kappa_\beta/16}
 \label{eq:finite-density-yamabe-transfer}
\end{gather}
whenever the displayed denominator is positive.  Indeed, writing
$e=s_g+4$, one has $\int e^2\le4\kappa_\beta V$ and
$16\widehat V=16V-8\int e+\int e^2-\delta$ with
$\delta\le D_\beta\le\kappa_\beta V$.  This gives
$\widehat V/V\ge1-\sqrt{\kappa_\beta}-\kappa_\beta/16$;
the nonincreasing numerator gives
\eqref{eq:finite-density-yamabe-transfer}.

The assertion at $\beta=2$ is the nonnegativity of the left side.
If no $c(B)>0$ worked for $\beta>2$, there would be admissible
metrics and parameters with
\[
 \frac{\beta_j^2\kappa_{\beta_j}(g_j)}{(\beta_j-2)^2}
       \longrightarrow0.
\]
In particular $\kappa_{\beta_j}(g_j)\to0$, and
\eqref{eq:finite-density-yamabe-transfer} yields the same vanishing
ratio for $\widehat\kappa_{\beta_j}$.  Since $a_j^2\ge T_{+,j}$,
Theorem~\ref{thm:single-parameter} gives
\[
 \widehat\kappa_{\beta_j}\ge
 \frac{\widehat P_j}{\widehat V_j}
       +\frac{\widehat D_{\beta_j}}{\beta_j\widehat V_j}
 \ge c_{\rm par}\frac{(\beta_j-2)^2}{\beta_j^4}.
\]
Thus $\beta_j\to\infty$; this conclusion also excludes a sequence
approaching $\beta=2$ from above.  The hatted sequence now satisfies
\eqref{eq:finite-density-zero-sequence}, with the same projection
budget, square condition, and Euler-density bound.  It has just been
excluded by \eqref{eq:finite-density-contradiction}.
\end{proof}

\subsection{Polynomial dependence and the remaining endpoint}

The logarithmic line argument proves positivity of $c(B)$.  A
polynomial lower bound uses the separate positive-potential fifth-line
test of Section~\ref{sec:web-positive-potential}, its rate-sensitive
extension in Section~\ref{sec:web-density-rate}, and the quantitative
transfer from the original defect in Section~\ref{sec:web-primitive-rate}.

\begin{corollary}[A polynomial density bound]
\label{cor:finite-density-polynomial}
Fix the absolute slow-truncation exponent $0<\epsilon_0\le1$ in
\eqref{eq:finite-density-slow-truncation}, and set
\[
 \gamma=\frac{\epsilon_0}{1+\epsilon_0},\qquad
 N_0=\frac{24}{\gamma}.
\]
There is an absolute $c_0>0$ such that the constant in
Theorem~\ref{thm:finite-density-endpoint} may be chosen with
\begin{equation}
 c(B)\ge c_0(1+B)^{-N_0}.
 \label{eq:finite-density-polynomial}
\end{equation}
On the scalar-minus-four representative, writing
$b=\fint|W^-|^2$ and $x=\kappa_\beta+\beta^{-2}$, one also has
\begin{equation}
 x\ge c_0(1+b)^{-N_0},
 \label{eq:finite-density-curvature-rate}
\end{equation}
after reducing the absolute constant.  No optimality of $N_0$ is
claimed.
\end{corollary}

\begin{proof}
Let $\mathfrak e$ be the fixed sum of primitive actual errors defined
in Section~\ref{sec:web-density-rate}.  It includes the two
amplitude errors, the resolvent and Dirac residuals, the full first
projection, both the $H^1$ and $L^4$ differences $v-t$, the fixed-cutoff
phase and connection-fitting errors, the determinant-curvature fit,
the Ricci and positive-curvature errors, and $\beta^{-1}$.
Section~\ref{sec:web-primitive-rate} proves, with constants independent
of $b$,
\begin{equation}
 \mathfrak e\le Cx^{\gamma/16}.
 \label{eq:finite-density-primitive-rate}
\end{equation}
The exponent includes the $L^4$ transfer loss.  It bounds the specified
measured sum, not an error sequence chosen after observing $b$.
Section~\ref{sec:web-density-rate} excludes an actual endpoint
sequence satisfying
\begin{equation}
 \mathfrak e_j(1+b_j)^{3/2}\longrightarrow0.
 \label{eq:finite-density-positive-potential-exclusion}
\end{equation}
These are the two additional inputs to this corollary.

If \eqref{eq:finite-density-curvature-rate} failed uniformly, a
sequence could be chosen with $x_j(1+b_j)^{N_0}\to0$.  Then
$x_j\to0$, and \eqref{eq:finite-density-primitive-rate} gives
\[
 \mathfrak e_j(1+b_j)^{3/2}
 \le C\bigl[x_j(1+b_j)^{24/\gamma}\bigr]^{\gamma/16}
       \longrightarrow0,
\]
contrary to \eqref{eq:finite-density-positive-potential-exclusion}.
This proves \eqref{eq:finite-density-curvature-rate}.

If $\kappa_\beta\le1$, the Euler estimate
\eqref{eq:finite-density-opposite-weyl} gives
$1+b\le C(1+B)$.  After reducing constants, choose
$m_B=c(1+B)^{-N_0}\le1/2$ so that
$\kappa_\beta+\beta^{-2}\ge m_B$.  Put
$\beta_B=(2/m_B)^{1/2}$.  For $\beta\ge\beta_B$,
$\kappa_\beta\ge m_B/2$.  For $2\le\beta\le\beta_B$,
the finite-parameter estimate gives
\[
 \kappa_\beta\ge c_{\rm par}\frac{(\beta-2)^2}{\beta^4}
       \ge\frac{c_{\rm par}m_B}{2}
                    \frac{(\beta-2)^2}{\beta^2}.
\]
The branch $\kappa_\beta>1$ is immediate after reducing the
constant.  Thus the desired polynomial inequality holds on
$s=-4$ representatives.  For an original metric with small
$\kappa_\beta$, \eqref{eq:finite-density-yamabe-transfer} gives
$\widehat\kappa_\beta\le2\kappa_\beta$; the complementary branch
again costs only an absolute constant.  This proves
\eqref{eq:finite-density-polynomial} for all metrics in the theorem.
\end{proof}

An unrestricted sequence with $\beta_j\to\infty$ and
$\kappa_{\beta_j}\to0$ must therefore satisfy
$\chi(M_j)/A_j\to\infty$; otherwise a bounded-density subsequence
would contradict Theorem~\ref{thm:finite-density-endpoint}.
Equation~\eqref{eq:finite-density-curvature-rate} further restricts
its possible growth rate, but permits sufficiently rapid opposite-Weyl
blow-up.  The finite-density proof uses $B$ in
\eqref{eq:finite-density-negative-hessian}, in the mass-gradient and
covariance tests, and in the $H^1$ bound for the fifth-line section.
These estimates have not been made uniform as $B\to\infty$.
The density-free endpoint, an endpoint theorem for arbitrary bounded
self-dual forms, and the balanced first-order coefficient remain
separate questions.
